\documentclass[10pt]{article}
\usepackage[margin=0.78in]{geometry}
\usepackage[T1]{fontenc}
\usepackage{lmodern}
\usepackage{microtype}
\usepackage{booktabs}
\usepackage{longtable}
\usepackage{tabularx}
\usepackage{enumitem}
\usepackage{hyperref}
\usepackage{xcolor}
\usepackage{listings}
\usepackage{amsmath,amssymb}

\hypersetup{colorlinks=true,linkcolor=blue!55!black,urlcolor=blue!55!black}
\setlist{nosep,leftmargin=*}
\lstset{basicstyle=\ttfamily\small,breaklines=true,columns=fullflexible,
  frame=single,showstringspaces=false}

\title{Plain2MeTTa General Semantic Validation\\
       Architecture and Integration Specification}
\author{Working specification for review}
\date{2026-08-16, revision 0.2}

\begin{document}
\maketitle

\begin{abstract}
This document specifies a general semantic-validation capability for
Plain2MeTTa.  The capability converts reviewed Plain requirements into typed
contracts, independently generated validation plans, executable oracles,
properties, and formal obligations.  Generated MeTTa and grounded Python are
then evaluated under bounded runtimes, with exact provenance from source text
to evidence.  The design does not claim that arbitrary English can be proved
correct automatically.  It makes the interpretation boundary explicit,
requires human approval where meaning is uncertain, grades the strength of
evidence, and fails closed rather than silently manufacturing semantics.
\end{abstract}

\tableofcontents

\section{Purpose and non-goals}

The current Plain2MeTTa evaluation workbench demonstrates three reviewed
examples that pass through Plain, a logical intermediate representation (IR),
generated MeTTa and Python, and bounded dual-runtime checks.  Those examples
use exact, curated behavioral profiles.  This specification generalizes that
mechanism without replacing it with an ever larger table of example-specific
code.

The intended result is a framework in which each requirement is associated
with explicit semantic obligations and independently inspectable evidence.
The framework must answer five different questions separately:

\begin{enumerate}
\item Was the source text preserved and reviewed?
\item Is the logical contract well formed and internally consistent?
\item Did generated artifacts execute in their declared runtimes?
\item Did observations satisfy independently specified examples and
      properties?
\item Is any stronger statement formally established, and under which
      assumptions?
\end{enumerate}

Non-goals include unrestricted execution of generated code, automatic proof
of arbitrary natural-language intent, treating model consensus as proof,
silently filling policy or domain gaps, and collapsing MeTTa reasoning into
Python implementation details.

\section{Normative terminology and evidence grades}

The words MUST, MUST NOT, SHOULD, SHOULD NOT, and MAY are normative.

\begin{longtable}{p{0.15\linewidth}p{0.77\linewidth}}
\toprule
Grade & Meaning \\
\midrule
G0: preserved & Exact reviewed source bytes, identities, and hashes are known. \\
G1: structured & The contract and validation plan pass schema, type, provenance,
and consistency checks. \\
G2: executable & The generated artifact starts and terminates under its declared
bounded runtime. \\
G3: examples & Independent concrete input/output examples pass. \\
G4: properties & Property-based, metamorphic, state-machine, or differential
checks pass over a recorded test domain. \\
G5: proved & A stated formula is proved or exhaustively checked in a named formal
system under explicit assumptions and bounds. \\
G6: observed & A deployed system satisfies monitored obligations over a recorded
operational interval; this is not a universal proof. \\
\bottomrule
\end{longtable}

Grades are not interchangeable.  G2 does not imply G3.  G4 does not imply G5.
A result MUST report a vector of grades per obligation rather than one global
``semantically valid'' Boolean.  A convenient summary MAY be the weakest grade
required by an approved acceptance policy.

\section{Trust model and independence}

Semantic validation is vulnerable to correlated mistakes.  If one model or
one deterministic transformation writes both the implementation and its
oracle, both can encode the same misunderstanding.  Plain2MeTTa therefore
separates four roles:

\begin{description}
\item[Interpretation author] Converts reviewed Plain into typed contracts and
marks unresolved meaning.
\item[Implementation author] Produces MeTTa and grounded modules from approved
contracts.
\item[Validation author] Produces examples, properties, generators, and oracle
logic from the approved source and contract, without reading generated
implementation bodies by default.
\item[Evidence runner] Executes immutable implementation and validation
artifacts under bounded, identified runtimes and records observations.
\end{description}

The roles MAY use different models, deterministic compilers, or humans.  Their
inputs, prompts, versions, and outputs MUST be recorded.  The validation author
MUST NOT derive expected output by invoking the implementation under test.
When full independence is impossible, the shared dependency MUST be named and
the evidence downgraded.

\section{Extension to the logical IR}

The existing logical IR already carries contracts, obligations, dependencies,
invariants, source clauses, and operational holes.  The following immutable,
content-addressed artifacts extend it.

\subsection{SemanticContract}

\begin{lstlisting}
SemanticContract := {
  contract_id: CanonicalID,
  source_clause_refs: [ArtifactSpanRef, ...],
  name: Identifier,
  inputs: [TypedVariable, ...],
  output: TypeExpression,
  preconditions: [Predicate, ...],
  postconditions: [Predicate, ...],
  invariants: [Predicate, ...],
  effects: [EffectDeclaration, ...],
  temporal_constraints: [TemporalFormula, ...],
  nondeterminism: DeterminismPolicy,
  environment_assumptions: [Assumption, ...],
  unresolved_holes: [HoleRef, ...]
}
\end{lstlisting}

Predicates MUST use a typed, finite core language.  An extension MAY embed
MeTTa terms, but it MUST declare evaluation semantics and purity.  Free prose
MAY accompany a predicate but MUST NOT substitute for an executable or formal
meaning.

\subsection{ValidationObligation}

\begin{lstlisting}
ValidationObligation := {
  obligation_id: CanonicalID,
  contract_ref: ArtifactRef,
  source_clause_refs: [ArtifactSpanRef, ...],
  claim: Predicate | TemporalFormula | StatisticalClaim,
  required_grade: EvidenceGrade,
  admissible_methods: [MethodKind, ...],
  domain: TestDomain,
  assumptions: [AssumptionRef, ...],
  severity: critical | major | minor,
  unresolved: bool
}
\end{lstlisting}

No obligation may be considered discharged if it is unresolved, references a
stale contract, or relies on an unapproved assumption.

\subsection{ValidationPlan}

\begin{lstlisting}
ValidationPlan := {
  plan_id: CanonicalID,
  reviewed_contract_refs: [ArtifactRef, ...],
  author_provenance: AdapterProvenance,
  examples: [ExampleOracle, ...],
  generators: [InputGenerator, ...],
  properties: [PropertyOracle, ...],
  metamorphic_relations: [MetamorphicRelation, ...],
  state_models: [StateMachineOracle, ...],
  differential_oracles: [DifferentialOracle, ...],
  formal_tasks: [FormalTask, ...],
  coverage_claims: [CoverageClaim, ...]
}
\end{lstlisting}

The plan is reviewed independently and approved against exact contract bytes.
Changing a source clause, contract, assumption, generator, or oracle MUST
invalidate downstream results.

\subsection{RuntimeEvidence and ValidationVerdict}

Runtime evidence records the runtime identity and hash, resource bounds, input
seed or case identity, complete normalized observations, exit status, and
artifact hashes.  A verdict joins evidence to one exact obligation and method:

\begin{lstlisting}
ValidationVerdict := {
  verdict_id: CanonicalID,
  obligation_ref: ArtifactRef,
  implementation_ref: ArtifactRef,
  validation_plan_ref: ArtifactRef,
  runtime_evidence_refs: [ArtifactRef, ...],
  grade_achieved: EvidenceGrade,
  status: pass | fail | unknown | blocked,
  counterexample_refs: [ArtifactRef, ...],
  residual_risk: Text,
  assumptions_used: [AssumptionRef, ...]
}
\end{lstlisting}

\section{Predicate and oracle language}

The first implementation SHOULD use a deliberately small contract language.
It needs scalar and algebraic data types, records, finite collections, pure
functions, equality and ordering, quantifiers over bounded domains, event
traces, and explicit time.  It MUST prohibit implicit I/O and unbounded
quantification in executable oracles.

Example core forms are:

\begin{lstlisting}
(requires contract-id predicate)
(ensures contract-id predicate)
(invariant contract-id predicate)
(always trace-predicate)
(before event-a event-b)
(exactly-once event-pattern)
(for-all (x generator-id) predicate)
(approximately expected observed tolerance confidence-policy)
\end{lstlisting}

Each form has both a canonical serialized representation and one or more
runtime backends.  A deterministic reference interpreter SHOULD define the
meaning of the core.  MeTTa MAY own symbolic reduction, contract composition,
and provenance; Python MAY ground numerical, cryptographic, file, and systems
operations.  The ownership locus MUST be declared per predicate.

\section{Validation methods}

\subsection{Concrete examples}

Examples bind explicit inputs, initial state, permitted environment events,
and expected observations.  They are useful but weak.  Boundary examples and
negative examples are mandatory for critical obligations.

\subsection{Property-based testing}

An input generator is itself a reviewed artifact with a finite budget,
distribution description, shrinker, and seed.  On failure, the smallest found
counterexample is preserved.  Passing $N$ generated cases MUST be reported as
bounded evidence, never as proof over an infinite domain.

\subsection{Metamorphic testing}

When an exact answer is difficult to compute, a transformation on input can
imply a relation on output.  Examples include permutation invariance,
monotonicity, scale equivariance, idempotence, round trips, conservation, and
irrelevance of data outside an authorized window.

\subsection{State-machine and trace testing}

Stateful services require commands, preconditions, transitions, observations,
and invariants.  Generated command sequences explore reachable states.  Trace
oracles validate ordering, cardinality, secrecy, and temporal properties.

\subsection{Differential and cross-runtime testing}

When MeTTa and Python implement the same declared observable contract, they
MUST be run on identical canonical inputs and their normalized outputs compared.
Agreement is useful but not sufficient: two implementations can share a bad
contract or generator.  A reference interpreter or independent oracle remains
necessary.

\subsection{Formal methods}

Finite domains and pure contract fragments MAY be translated to SAT/SMT,
model checking, or exhaustive MeTTa reduction.  Every proof result MUST include
the exact theorem, assumptions, solver/runtime identity, resource bounds, and
proof certificate or replay command when available.  Unknown and timeout are
not pass.

\subsection{Mutation and adversarial validation}

The framework SHOULD mutate generated implementations and contracts.  Useful
mutations include inverted comparisons, removed state updates, wrong event
order, boundary shifts, leaked fields, changed time windows, and duplicated
effects.  A validation plan that fails to reject relevant mutants has weak
discriminative power and MUST expose that limitation.

\section{Pipeline integration}

The proposed extension preserves the current immutable phase chain.

\begin{enumerate}
\item \textbf{Source and elaboration.} Preserve original Plain bytes.  Mark
every added interpretation as derived, guidance-specified, advisor-proposed,
or unresolved.
\item \textbf{Human review.} Review the elaborated specification and tests.
Edits create new immutable snapshots.
\item \textbf{Contract IR.} Produce typed contracts, assumptions, effects,
and unresolved holes from reviewed snapshots.
\item \textbf{Logical review.} Detect inconsistent types, contradictory
invariants, unreachable obligations, temporal cycles, leakage, and missing
groundings.  Critical open findings block later claims.
\item \textbf{Independent validation synthesis.} Generate a ValidationPlan
using a separate adapter boundary.  The adapter sees reviewed source and
contracts, but not generated implementation bodies.
\item \textbf{Plan review and approval.} A person or policy approves exact
validation-plan bytes and any risk acceptance.
\item \textbf{Compilation.} Generate MeTTa and grounded modules from approved
contracts.  Record compiler provenance and exact generated-file hashes.
\item \textbf{Bounded execution.} Run implementation and oracle artifacts in
separate sandboxes where practical.  Network is denied by default; secrets are
absent; file and process access are bounded.
\item \textbf{Verdict.} Join exact obligations, plans, generated artifacts,
runtimes, observations, and counterexamples into graded verdicts.
\item \textbf{Traceability.} Reconstruct the full chain on every query.  Any
stale or mismatched reference fails closed.
\end{enumerate}

\section{Service and Web UI requirements}

The Web UI SHOULD show a left-to-right evidence graph:

\begin{center}
Plain $\rightarrow$ reviewed text $\rightarrow$ contracts $\rightarrow$
validation plan $\rightarrow$ code $\rightarrow$ evidence $\rightarrow$ verdict
\end{center}

For each obligation it MUST show source text, interpretation status, contract,
required evidence grade, plan author, implementation author, runtime identity,
cases or proof tasks, counterexamples, achieved grade, and residual risk.
Generated, executed, tested, proved, and operationally observed labels MUST be
visually distinct.

Suggested narrow APIs are:

\begin{lstlisting}
POST /api/validation-plan/<project-id>
GET  /api/validation-plan/<project-id>
POST /api/validation-plan-review/<project-id>
POST /api/semantic-test/<project-id>
GET  /api/semantic-results/<project-id>
GET  /api/semantic-trace/<project-id>?obligation=<canonical-id>
\end{lstlisting}

POST routes MUST use bounded duplicate-safe JSON, exact canonical identities,
explicitly injected adapters, and atomic admission.  They MUST NOT accept
provider credentials or arbitrary commands.  GET routes return metadata by
default; potentially sensitive bodies require separate authorization.

\section{Qualitatively different examples}

\subsection{Example A: pure numerical transformation}

\paragraph{Plain.}
``Given real $x$, lower bound $a$, and upper bound $b$ with $a<b$, return
$(x-a)/(b-a)$ clipped to $[0,1]$.  The function must be monotone in $x$.''

\paragraph{Contract.}
Inputs are finite reals; precondition $a<b$; output $y$ is finite and
$0\le y\le1$; $x\le x'$ implies $f(x,a,b)\le f(x',a,b)$; $f(a,a,b)=0$ and
$f(b,a,b)=1$.

\paragraph{Validation.}
Concrete boundary examples, generated finite triples, monotonic pair tests,
translation metamorphism $f(x+c,a+c,b+c)=f(x,a,b)$, differential MeTTa/Python
comparison, and an SMT proof over an exact rational implementation.  Mutants
that omit clipping or reverse subtraction must be killed.

\subsection{Example B: stateful authentication service}

\paragraph{Plain.}
``A successful password reset revokes all existing sessions before success is
reported.  Failed login and locked-account responses are indistinguishable to
the client.  Passwords and bearer tokens never appear in logs.''

\paragraph{Contract.}
The state contains users, password verifiers, sessions, and an append-only
audit trace.  Reset success requires every prior session for the user to be
revoked in an earlier trace position.  Public failure observations are equal
for bad credentials and lockout.  A secrecy predicate rejects credential
substrings and structured secret fields in every emitted log event.

\paragraph{Validation.}
A model-based state-machine runner generates login, refresh, reset, revoke,
and inspect-log commands.  Trace properties check ordering and noninterference.
Metamorphic tests replace one invalid password with another and require the
same public observation.  Fault injection interrupts the reset transaction.
The validator must reject implementations that report success before
revocation, leak a token, or distinguish lockout.

\subsection{Example C: time-series model without future leakage}

\paragraph{Plain.}
``At prediction time $t$, features use no observation after $t$.  Normalization
is fit on training data only.  Hyperparameters use validation data, never test
data.  The selected model is evaluated exactly once on the test partition.''

\paragraph{Contract.}
Every feature value carries an event-time lineage set.  The maximum lineage
time is at most its prediction time.  Scaler-fit lineage is a subset of the
training partition.  Selection traces contain no test identifiers.  The test
evaluation event has cardinality one and occurs after selection is frozen.

\paragraph{Validation.}
Synthetic datasets include a future-only sentinel.  Perturbing all future
observations must not affect earlier predictions.  Shuffling or changing test
labels must not change selected hyperparameters.  Provenance inspection checks
lineage.  Mutation tests introduce a centered rolling window, global scaler,
or repeated test evaluation and require counterexamples.  Statistical metric
thresholds are separate claims with confidence and dataset assumptions; they
must not be confused with leakage invariants.

\subsection{Example D: idempotent distributed command handling}

\paragraph{Plain.}
``A transfer command has a client request ID.  Retrying the same command may
return the prior result but must not debit twice.  Two distinct request IDs may
represent distinct transfers.  A crash between durable debit and response must
remain recoverable.''

\paragraph{Contract.}
The observable state includes balances, a durable request ledger, and emitted
responses.  For one request ID and payload, at most one debit transition is
committed.  Reuse of an ID with a different payload is rejected.  After any
allowed crash point, replay converges to a state observationally equivalent to
one uninterrupted execution.

\paragraph{Validation.}
The state-machine generator produces retries, reorderings, duplicate delivery,
and crashes at named persistence boundaries.  Metamorphic replay duplicates
messages without changing final balances.  A small-state model checker explores
two accounts, bounded amounts, and finite crash points.  The runtime harness
injects failures and compares MeTTa orchestration with a Python grounded store.
Mutants that record the request after debit or key only by payload must fail.

\subsection{Example E: policy-dependent data export}

\paragraph{Plain.}
``Support staff may export retained audit events only when authorized by the
current policy.''

This example is intentionally incomplete.  ``Current policy'', retention,
staff identity, authorization source, and permitted fields are undefined.  The
correct general validator creates typed unresolved holes and review questions.
It may validate mechanics such as ``no export occurs without an authorization
decision'', but MUST NOT invent a retention period or role policy.  Compilation
or a strong verdict remains blocked until the holes are resolved and approved.

\section{Acceptance criteria for the first general release}

The first release is accepted only if all of the following hold:

\begin{enumerate}
\item At least four qualitatively different examples use data-driven contracts
and plans, not code branches keyed to exact source text.
\item At least one example exercises pure properties, one state-machine trace,
one temporal/data-lineage invariant, and one unresolved policy hole.
\item MeTTa and Python execute through pinned bounded runtimes where both are
applicable, and cross-runtime comparison is explicit.
\item Property generation records seeds and shrinking; counterexamples are
immutable artifacts linked to source obligations.
\item A meaningful mutation suite demonstrates that validation rejects seeded
semantic defects.

\section{Recommended software toolchain}

No single formal tool is appropriate for every obligation.  Plain2MeTTa SHOULD
route obligations by semantic shape and required assurance, while preserving a
common immutable evidence envelope.

\subsection{Hypothesis for executable properties and state machines}

Python Hypothesis SHOULD be the default engine for generated examples,
property-based tests, metamorphic relations, and executable state-machine
testing.  In particular, \texttt{RuleBasedStateMachine} supports commands,
preconditions, invariants, generated command sequences, reproducible seeds,
and automatic counterexample shrinking.

The Plain2MeTTa adapter MUST generate ordinary reviewable Python test modules,
not hidden dynamic calls.  Each run records the Hypothesis version, settings,
seed or database identity, maximum examples, deadlines, phases, and shrunk
counterexample.  Health-check suppression MUST be explicit and justified.
Hypothesis is the fast default bug-finding layer; passing a finite campaign is
G4 bounded evidence, not a proof.

\subsection{TLA+ and TLC for temporal and distributed state machines}

TLA+ SHOULD be the canonical model language for concurrency, temporal
ordering, retry/idempotency protocols, crash recovery, fairness, and safety or
liveness invariants.  TLC SHOULD be the initial model checker because it has a
mature semantics, explicit finite configurations, and concrete counterexample
traces.

A generated TLA+ bundle consists of a module, finite model configuration,
source-to-variable map, invariant and temporal-property map, and normalization
rules for counterexample traces.  A TLC pass is valid only for the exact finite
scope recorded in the evidence.  State-space size, distinct states, depth,
workers, symmetry sets, and liveness mode MUST be recorded.  Deadlock checking
MUST NOT be disabled silently.

Apalache MAY be added as a symbolic backend for bounded checking when explicit
TLC exploration becomes too large.  Quint MAY be offered later as a friendlier
authoring or simulation surface, but revision 0.2 keeps generated TLA+ as the
canonical temporal artifact to avoid an unnecessary semantic translation in
the first implementation.

\subsection{Z3 through SMT-LIB for bounded logical obligations}

Z3 SHOULD discharge satisfiability, finite arithmetic, reachability, type and
contract consistency, bounded transition, and counterexample-generation tasks.
Plain2MeTTa MUST emit canonical SMT-LIB text plus a declaration/source map,
rather than persist opaque solver API objects.  The evidence records Z3
version, command, logic, options, timeout, result, model or unsat core, and
proof object when requested and supported.

The trusted claim is exactly the encoded formula.  Therefore an independently
tested encoder and round-trip fixtures are mandatory.  \texttt{unknown},
timeout, resource exhaustion, or unsupported theory yields Unknown, never Pass.
Solver satisfiability does not establish that the original English was
interpreted correctly.

\subsection{Lean 4 and Mathlib for durable theorem-level assurance}

Lean 4 with Mathlib SHOULD define the small trusted Plain2MeTTa contract
calculus and its denotational or operational semantics.  Lean is the preferred
G5 layer for stable claims such as:

\begin{itemize}
\item well-typed contracts do not produce malformed observations;
\item canonical serialization and IR normalization preserve meaning;
\item approved lowering rules refine the contract semantics;
\item selected pure algorithms satisfy their contracts;
\item a finite-state transition model preserves a named invariant; and
\item transformations used by test generation are sound for their stated
metamorphic relation.
\end{itemize}

Lean SHOULD NOT be required for every generated runtime test.  Proof cost is
reserved for the semantic kernel, reusable lowering theorems, high-risk pure
algorithms, and explicitly requested application claims.  Each Lean artifact
pins the Lean toolchain, Mathlib revision, imports, theorem statement, axioms,
and build command.  Use of \texttt{sorry}, new axioms, unsafe declarations, or
unreviewed native decision procedures MUST prevent a proved verdict.

Where Z3 or TLC discovers a bounded result, Lean MAY separately prove a
general theorem, but a solver's success MUST NOT be relabeled as a Lean proof.
Likewise, generated proof scripts require kernel checking before admission.

\subsection{Hyperon and Python runtime conformance}

The existing pinned Hyperon CLI and bounded Python subprocess remain the
implementation-evidence layer.  Both receive identical canonical cases where
their observable contracts overlap.  Normalized output is compared with the
independent oracle and with the other runtime.  Cross-runtime agreement raises
confidence but does not replace Hypothesis, TLA+, SMT, or Lean evidence.

\subsection{Routing matrix}

\begin{longtable}{p{0.24\linewidth}p{0.19\linewidth}p{0.47\linewidth}}
\toprule
Obligation shape & Primary tool & Secondary evidence \\
\midrule
Pure executable function & Hypothesis & Z3 for bounded arithmetic; Lean for
selected general theorem \\
Stateful sequential API & Hypothesis state machine & Lean for reusable state
invariant if warranted \\
Concurrent/distributed protocol & TLA+/TLC & Hypothesis fault-injection;
Apalache for symbolic bounds \\
Contract consistency & Z3/SMT-LIB & deterministic IR checks; Lean kernel
metatheory \\
Temporal/data-lineage rule & TLA+ or SMT by shape & Hypothesis metamorphic
tests and runtime provenance \\
Compiler/IR semantic preservation & Lean 4 & differential fixtures and
mutation testing \\
Concrete MeTTa/Python behavior & Hyperon/Python runtimes & independent oracle,
Hypothesis, and differential comparison \\
Unresolved policy or world fact & no automatic tool & explicit hole and human
review \\
\bottomrule
\end{longtable}
\item Changing source, contracts, assumptions, plans, generated code, or
runtimes invalidates downstream verdicts.
\item The UI displays evidence grades and never labels unsupported arbitrary
input as semantically validated.
\item Focused tests, the complete provider-free suite, live browser/API smoke,
runtime smoke, repository hygiene, and credential scans pass.
\end{enumerate}

\section{Step-by-step plan for coding agents}

Each stage is a separately reviewable task-branch increment.  Agents MUST read
the current project records, repository instructions, CI, and applicable
skills before editing.  They MUST preserve unrelated work, use immutable
artifacts and exact hashes, run focused tests before the full suite, and create
a reproducible experiment record.  An agent MUST NOT begin a dependent stage
until the preceding acceptance gate passes.

\subsection*{Stage 0: freeze interfaces and threat model}

\textbf{Work.} Inventory current logical-IR, compiler-output, sandbox,
traceability, evaluation UI, and dual-runtime schemas.  Write the canonical
evidence-grade vocabulary, role-separation policy, runtime trust boundaries,
and migration compatibility rules.  Select and pin candidate versions of
Hypothesis, TLC, Z3, Lean 4/Mathlib, and Hyperon in isolated environments.

\textbf{Gate.} A source-located interface map exists; upstream smoke tests for
each chosen tool pass unchanged; no system-wide dependency change occurs; and
the existing 662-test baseline and live three-example evaluation remain green.

\subsection*{Stage 1: immutable semantic artifact model}

\textbf{Work.} Implement strict schemas and canonical serialization for
\texttt{SemanticContract}, \texttt{ValidationObligation},
\texttt{ValidationPlan}, generators, oracles, runtime evidence,
counterexamples, and graded verdicts.  Bind every artifact to exact reviewed
source and contract references.  Extend transitive invalidation.

\textbf{Gate.} Round-trip and adversarial-deserialization tests pass; forged,
partial, duplicate, stale, path-confused, and hash-mismatched artifacts fail
closed; changing any upstream byte invalidates every dependent verdict.

\subsection*{Stage 2: finite contract calculus and interpreter}

\textbf{Work.} Define the typed predicate core, effects, bounded quantifiers,
trace primitives, time semantics, assumption references, and ownership locus.
Implement a deterministic reference interpreter and canonical MeTTa projection.
Keep unsupported constructs as typed holes.

\textbf{Gate.} Gold fixtures cover pure, stateful, temporal, approximate, and
unresolved contracts.  Type preservation and interpreter determinism pass
provider-free tests.  No free prose is executed as a predicate.

\subsection*{Stage 3: independent validation-plan synthesis}

\textbf{Work.} Add a one-call, provider-independent validation-author adapter
that sees reviewed source and contracts but not generated implementation
bodies.  Admit only strict plans.  Add plan review, edits, decisions, and exact
approval binding.

\textbf{Gate.} Role separation is demonstrated by tests; malformed or
misattributed adapter output causes no partial writes or retry; unresolved
critical meaning blocks approval; source or contract edits invalidate the plan.

\subsection*{Stage 4: Hypothesis backend}

\textbf{Work.} Lower examples, generators, pure properties, metamorphic
relations, and sequential state machines to reviewable Hypothesis modules.
Run them in the bounded Python sandbox.  Persist settings, seeds, cases,
shrinking history, and minimal counterexamples.

\textbf{Gate.} Numerical, authentication, and idempotency fixtures pass.
Seeded mutants yield stable shrunk counterexamples.  A zero-exit run with a
wrong observation fails.  Replaying recorded seeds reproduces results.

\subsection*{Stage 5: TLA+/TLC backend}

\textbf{Work.} Lower finite transition contracts to TLA+ modules and TLC model
configurations with source maps.  Parse TLC states and counterexample traces
into bounded immutable evidence.  Initially cover authentication ordering and
distributed idempotency/crash recovery.

\textbf{Gate.} Correct models satisfy declared safety properties in recorded
finite scopes; injected ordering, duplicate-debit, and recovery mutants produce
source-linked counterexample traces.  State counts, depth, scope, fairness, and
deadlock configuration are visible in the UI.

\subsection*{Stage 6: SMT-LIB/Z3 backend}

\textbf{Work.} Lower the supported pure and bounded-transition fragment to
canonical SMT-LIB.  Add exact declaration maps, model/unsat-core parsing,
timeouts, and Unknown handling.  Independently test the encoder using gold
formulas and mutation pairs.

\textbf{Gate.} Contradictory contracts, unreachable states, boundary errors,
and finite counterexamples are detected.  Unknown and timeout fail closed.
No solver result is admitted without exact formula and solver provenance.

\subsection*{Stage 7: Lean 4 semantic kernel}

\textbf{Work.} Create a small Lean package defining types, values, predicates,
contract satisfaction, traces, and the supported lowering relation.  Prove
core well-formedness and selected preservation theorems.  Add one pure example
proof and one finite state-invariant proof.  Export theorem identities and
source maps as evidence metadata.

\textbf{Gate.} Pinned \texttt{lake build} passes with no \texttt{sorry}, new
axioms, unsafe declarations, or unrecorded native trust.  Deliberately false
theorems fail.  The UI displays exact theorem statements and assumptions and
does not describe TLC/Z3 results as Lean proofs.

\subsection*{Stage 8: dual-runtime and cross-tool verdict composition}

\textbf{Work.} Generalize the existing Hyperon/Python runner to consume
approved contracts and plans rather than exact-text profiles.  Normalize
observations, compare both runtimes with independent oracles, and compose
Hypothesis, TLC, Z3, and Lean evidence into per-obligation grade vectors.

\textbf{Gate.} No single tool result can over-promote an obligation.  Shared
dependencies and assumptions are reported.  Wrong-output, stale-evidence,
cross-runtime divergence, and missing-proof cases produce the correct Fail,
Unknown, or Blocked status.

\subsection*{Stage 9: API and Web UI integration}

\textbf{Work.} Add the narrow plan, review, execution, result, and trace routes.
Extend the browser workbench with an evidence graph, tool-specific details,
counterexample views, assumption/hole review, and artifact downloads.

\textbf{Gate.} Browser and API tests cover exact identities, duplicate-safe
framing, bounds, stale chains, and authorization.  The UI distinguishes G0--G6
and never labels arbitrary unsupported input as semantically validated.

\subsection*{Stage 10: vertical acceptance suite and mutation score}

\textbf{Work.} Convert the current three profiles from source-text branches to
data artifacts.  Implement the five examples in Section 9, including the
intentionally blocked policy case.  Build a categorized mutation corpus.

\textbf{Gate.} At least four semantic shapes use the general pipeline; relevant
mutants are killed at an agreed threshold with surviving mutants reviewed;
focused and full tests, live dual-runtime replay, model checking, SMT checks,
Lean replay, secret/large-file/untracked audits, and exact trace reconstruction
all pass.  Publish only to a task branch and draft PR.

\subsection*{Stage 11: independent review and stabilization}

\textbf{Work.} Ben evaluates concrete UI examples.  Revise behavior before
writing a polished design report.  Then solicit frontier-model review with
frozen artifacts, incorporate accepted findings, and finally seek human MeTTa
and formal-methods developer review.

\textbf{Gate.} Review findings have explicit accept/reject/defer decisions and
regression evidence.  Default-branch merge or release remains a separate human
decision.

\section{Open design questions}

\begin{itemize}
\item Which contract core should be canonical: a restricted MeTTa vocabulary,
JSON with MeTTa projection, or both bound by a canonical hash?
\item Which property-testing engine should ground generation and shrinking,
and how should generators be represented for replay in MeTTa?
\item Which formal backends provide useful proof certificates without making
the core dependent on a large solver stack?
\item What separation between interpretation and validation authors is enough
for each risk tier?
\item Which claims require human approval, two-model review, mutation score,
formal proof, or operational monitoring?
\item How should probabilistic and approximate requirements express tolerances,
confidence, multiple comparisons, and distribution-shift assumptions?
\end{itemize}

\section{Summary}

General semantic validation in Plain2MeTTa should be an evidence-producing
contract system.  Plain meaning is reviewed into typed contracts; validation
plans are independently authored; MeTTa and grounded modules execute under
bounded runtimes; examples, properties, traces, mutations, and formal tasks
produce graded evidence; and every claim remains hash-bound to its exact source
and assumptions.  The architecture increases confidence without pretending
that execution, model agreement, or a finite test suite proves arbitrary
natural-language intent.

\end{document}
